% ==============================================================
%  Naawbi LMS --- CS3009 Software Engineering (Spring 2026)
%  FINAL REPORT  |  Team Ibwaan
%  Combined documentation across Sprints 1, 2, 3 (D1 + D2 + D3)
% ==============================================================
\documentclass[12pt, a4paper]{article}

% ── Packages ──────────────────────────────────────────────────
\usepackage[T1]{fontenc}
\usepackage[utf8]{inputenc}
\usepackage{lmodern}
\usepackage[top=2.3cm, bottom=2.3cm, left=2.3cm, right=2.3cm, headheight=28pt]{geometry}
\usepackage{microtype}
\usepackage{setspace}
\usepackage{parskip}
\usepackage{titlesec}
\usepackage{titling}
\usepackage{fancyhdr}
\usepackage{graphicx}
\usepackage{booktabs}
\usepackage{longtable}
\usepackage{multicol}
\usepackage{enumitem}
\usepackage{mdframed}
\usepackage{xcolor}
\usepackage{hyperref}
\usepackage{marvosym}
\newcommand{\faGithub}{\raisebox{-0.1ex}{\Keyboard}}
\usepackage{soul}
\usepackage{float}
\usepackage{caption}
\usepackage{subcaption}
\usepackage[most]{tcolorbox}

% ── Colours ───────────────────────────────────────────────────
\definecolor{primary}{HTML}{1A237E}
\definecolor{accent}{HTML}{283593}
\definecolor{light}{HTML}{E8EAF6}
\definecolor{rule}{HTML}{5C6BC0}
\definecolor{muted}{HTML}{546E7A}
\definecolor{rowalt}{HTML}{F3F4FA}
\definecolor{boxbg}{HTML}{EEF0FB}
\definecolor{success}{HTML}{1B5E20}
\definecolor{warn}{HTML}{B26A00}
\definecolor{danger}{HTML}{B71C1C}

% ── Hyperref ─────────────────────────────────────────────────
\hypersetup{
    colorlinks   = true,
    linkcolor    = primary,
    urlcolor     = accent,
    citecolor    = primary,
    pdftitle     = {Naawbi LMS — Final Report},
    pdfauthor    = {Team Ibwaan},
    pdfsubject   = {CS3009 Software Engineering Spring 2026},
    pdfkeywords  = {LMS, Software Engineering, Naawbi, Ibwaan, Final Report},
}

% ── Section formatting ───────────────────────────────────────
\titleformat{\section}
  {\Large\bfseries\color{primary}}
  {\thesection}{1em}{}
  [\color{rule}\titlerule]

\titleformat{\subsection}
  {\large\bfseries\color{accent}}
  {\thesubsection}{1em}{}

\titleformat{\subsubsection}
  {\normalsize\bfseries\color{accent}}
  {\thesubsubsection}{1em}{}

\titlespacing*{\section}{0pt}{1.6ex plus .5ex minus .2ex}{0.8ex}
\titlespacing*{\subsection}{0pt}{1.0ex plus .3ex minus .1ex}{0.4ex}

% ── Header / Footer ──────────────────────────────────────────
\pagestyle{fancy}
\fancyhf{}
\renewcommand{\headrulewidth}{0.4pt}
\renewcommand{\headrule}{\hbox to\headwidth{\color{rule}\leaders\hrule height \headrulewidth\hfill}}
\fancyhead[L]{\small\color{muted}\textbf{Naawbi LMS} — Final Report}
\fancyhead[R]{\small\color{muted}CS3009 Spring 2026}
\fancyfoot[C]{\small\color{muted}\thepage}

% ── Lists ─────────────────────────────────────────────────────
\setlist[itemize]{leftmargin=1.5em, itemsep=2pt, topsep=4pt}
\setlist[enumerate]{leftmargin=1.5em, itemsep=2pt, topsep=4pt}

% ── Image fallback ───────────────────────────────────────────
\newcommand{\imgOrPlaceholder}[3][width=0.85\textwidth]{%
  \IfFileExists{#2}{%
    \includegraphics[#1]{#2}%
  }{%
    \begin{center}%
    \fboxsep=8pt\fbox{\parbox{0.7\textwidth}{\centering\color{muted}%
      \textit{[Placeholder — image not available]}\\[4pt]%
      \footnotesize #3}}%
    \end{center}%
  }%
}

% ── Custom environments ───────────────────────────────────────
\tcbuselibrary{skins, breakable}

\newtcolorbox{infobox}[1][]{
  enhanced, breakable,
  colback=boxbg, colframe=rule,
  arc=4pt, boxrule=0.8pt,
  leftrule=4pt,
  title=#1,
  fonttitle=\bfseries\color{primary},
  #1
}

\newtcolorbox{userstory}[1][]{
  enhanced, breakable,
  colback=white, colframe=rule!40,
  arc=3pt, boxrule=0.5pt,
  top=4pt, bottom=4pt, left=8pt, right=8pt,
  #1
}

\newcommand{\fcol}[1]{\textbf{\color{primary}{#1}}}
\newcommand{\fcolplain}[1]{\textbf{#1}}

\begin{document}
\onehalfspacing

% ═══════════════════════════════════════════════════════════════
%  TITLE PAGE
% ═══════════════════════════════════════════════════════════════
\begin{titlepage}
  \centering

  \vspace*{2.0cm}

  {\small\scshape\color{muted} National University of Computer and Emerging Sciences}\\[0.4em]
  {\small\color{muted} CS3009 --- Software Engineering \quad|\quad Spring 2026}

  \vspace{1.2cm}
  {\color{rule}\hrule height 1.5pt}
  \vspace{1.2cm}

  {\fontsize{56}{62}\selectfont\bfseries\color{primary} Naawbi}\\[0.6em]
  {\Large\color{muted} Learning Management System}

  \vspace{1.0cm}
  {\color{rule!40}\rule{5cm}{0.6pt}}
  \vspace{1.0cm}

  {\LARGE\bfseries\color{primary} Final Report}\\[0.4em]
  {\normalsize\color{muted} Combined Project Documentation \,\textbar\, Sprints 1--3}

  \vspace{2cm}

  \begin{tabular}{r@{\quad}l}
    {\color{muted}\small Team Name}       & {\bfseries Ibwaan} \\[5pt]
    {\color{muted}\small Team Lead}       & {\bfseries Ibraheem Farooq} \\[5pt]
    {\color{muted}\small Members}         & 23i-0708 \textbar\ Ali Aan Khowaja \\
                                          & 23i-0816 \textbar\ Ibraheem Farooq \\
                                          & 23i-0819 \textbar\ Wajih-Ur-Raza Asif \\[5pt]
    {\color{muted}\small Submission Date} & 29 April 2026 \\[5pt]
    {\color{muted}\small Repository}      & \href{https://github.com/aliaankhowaja/Naawbi}{github.com/aliaankhowaja/Naawbi} \\
  \end{tabular}

  \vfill
  {\color{rule!40}\hrule height 0.5pt}
  \vspace{0.6cm}
  {\footnotesize\color{muted}
    \href{https://github.com/ibbi1020}{ibbi1020}\enspace\textbullet\enspace
    \href{https://github.com/wajih-rathore}{wajih-rathore}\enspace\textbullet\enspace
    \href{https://github.com/aliaankhowaja}{aliaankhowaja}
  }
  \vspace{0.6cm}
\end{titlepage}

% ── Table of Contents ─────────────────────────────────────────
\tableofcontents
\newpage

% ═══════════════════════════════════════════════════════════════
%  EXECUTIVE SUMMARY
% ═══════════════════════════════════════════════════════════════
\section*{Executive Summary}
\addcontentsline{toc}{section}{Executive Summary}

Naawbi is a desktop Learning Management System built by Team Ibwaan over three
sprints between February and April~2026 for the CS3009 Software Engineering
course. It is a single JavaFX~21 application backed by PostgreSQL~16, designed
for department-scale academic use --- a few courses with up to roughly 200
students each. The system delivers the complete academic workflow inside one
tool: instructors register, create courses, post announcements with file or
link attachments, create assignments with deadlines and late policies, grade
submissions, and view a class-wide gradebook; students register, enrol in
courses with a unique class code, submit work, see live submission status,
view per-course assignments and rosters, and track all pending work via a
deadline-aware To-Do list and a personal gradebook. Role-based access control
is enforced at every protected action so students never reach instructor-only
flows. Authentication uses SHA-256 password hashing and a session singleton
that is explicitly cleared on logout.

\medskip
\noindent\textbf{Status at submission.}
All 26 user stories scoped across the project (24 from the original D1
backlog plus the Assignments tab and People tab added during Sprint~3) have
been delivered and verified. Twenty-five formal black-box test cases were
executed against the seeded test database and all passed; both critical
controllers (\texttt{LoginController.handleLogin} and
\texttt{GradingController.handleSave}) reached 100\% branch coverage. Four
bugs were found during testing and all four were fixed before submission.

% ═══════════════════════════════════════════════════════════════
%  SECTION 1 — INTRODUCTION
% ═══════════════════════════════════════════════════════════════
\section{Introduction}

\subsection{Context}
This project is a Software Engineering team project to build a Learning
Management System (LMS) end-to-end --- requirements, design, implementation,
testing, and documentation --- and to deliver a working product covering the
full academic workflow for a course.

\subsection{Problem Statement}
Course management for small-to-mid scale classes typically suffers from
fragmented tools: one place for announcements, another for content, another
for submissions, another for grading, and yet another for tracking deadlines.
Students lack a clear, data-driven view of their performance and pending work.
Instructors need a single workflow to publish content, collect submissions,
grade, and monitor class progress. Naawbi addresses this fragmentation by
providing a single application that covers content, assignments, submissions,
grading, and a deadline-aware student dashboard.

\subsection{Goals}
\begin{itemize}
  \item Deliver a usable LMS that supports the end-to-end course workflow:
        \textbf{content} $\rightarrow$ \textbf{assignments} $\rightarrow$
        \textbf{submissions} $\rightarrow$ \textbf{grading}.
  \item Replace external deadline trackers with a built-in, urgency-coded
        To-Do list that aggregates pending work across all enrolled courses.
  \item Provide instructors with both a per-assignment grading interface and
        a class-wide gradebook macro-view, in one place.
\end{itemize}

% ═══════════════════════════════════════════════════════════════
%  SECTION 2 — PROJECT VISION
% ═══════════════════════════════════════════════════════════════
\section{Project Vision}

Build a clean, reliable, course-centric LMS that instructors use to manage
content and evaluations, and that students use to track progress and
performance, without needing any external tool for the core academic workflow.

\medskip
\noindent The two key differentiators are:
\begin{itemize}
  \item \textbf{Submission tracking with status derivation.} Every submission
        is timestamped server-side and automatically tagged \emph{Submitted /
        Late / Missing / Not Submitted} based on assignment deadline and
        instructor late policy.
  \item \textbf{Aggregated deadline awareness.} A student-facing To-Do list
        consolidates pending work across every enrolled course, with urgency
        colour-coding so overdue items and items due within 48~hours are
        immediately visible.
\end{itemize}

% ═══════════════════════════════════════════════════════════════
%  SECTION 3 — TEAM
% ═══════════════════════════════════════════════════════════════
\section{Team}

\begin{minipage}[c]{0.65\linewidth}
  \begin{tabular}{r@{\quad}l}
    \textbf{Team Name}    & Ibwaan \\
    \textbf{Project Name} & Naawbi \\
    \textbf{Team Lead}    & Ibraheem Farooq \\
  \end{tabular}
\end{minipage}%
\begin{minipage}[c]{0.3\linewidth}
  \raggedleft
  \imgOrPlaceholder[height=2cm]{../D3/assets/ibwaan_team_logo.png}{Team logo}
\end{minipage}

\vspace{0.5cm}

\subsection{Team Roles}

\begin{center}
\begin{tabular}{@{}p{0.28\linewidth}p{0.34\linewidth}p{0.32\linewidth}@{}}
  \toprule
  \textbf{Member} & \textbf{Role(s)} & \textbf{GitHub} \\
  \midrule
  Ibraheem Farooq    & PM, Scrum Master, Developer              & \href{https://github.com/ibbi1020}{github.com/ibbi1020} \\[4pt]
  Wajih-Ur-Raza Asif & Requirements Architect, Tester, Developer & \href{https://github.com/wajih-rathore}{github.com/wajih-rathore} \\[4pt]
  Ali Aan Khowaja    & Lead Developer, UI Designer              & \href{https://github.com/aliaankhowaja}{github.com/aliaankhowaja} \\
  \bottomrule
\end{tabular}
\end{center}

\subsection{Member Biographies}

\medskip

\noindent
\begin{minipage}[t]{0.30\linewidth}
  \vspace{0pt}
  \imgOrPlaceholder[width=\linewidth]{../D3/assets/ibraheem_photo.png}{Photo}
\end{minipage}%
\hspace{1em}%
\begin{minipage}[t]{0.66\linewidth}
  \vspace{0pt}
  {\bfseries\color{primary} Ibraheem Farooq}\\[-2pt]
  {\small\color{muted} PM / Scrum Master / Developer}\\[4pt]
  A Computer Science major based in Islamabad who transitioned from designing
  posters and brand identities to UX and product design. His technical background
  provides a unique perspective on bridging digital experiences that balance user
  needs with technical viability. Passionate about solving real problems through
  thoughtful design and clean, intuitive interfaces.
\end{minipage}

\medskip\noindent\color{rule!25}\hrule\color{black}\medskip

\noindent
\begin{minipage}[t]{0.30\linewidth}
  \vspace{0pt}
  \imgOrPlaceholder[width=\linewidth]{../D3/assets/wajih_photo.jpeg}{Photo}
\end{minipage}%
\hspace{1em}%
\begin{minipage}[t]{0.66\linewidth}
  \vspace{0pt}
  {\bfseries\color{primary} Wajih-Ur-Raza Asif}\\[-2pt]
  {\small\color{muted} Requirements Architect / Tester / Developer}\\[4pt]
  A software engineer focused on building practical systems that solve real
  operational problems. His background spans end-to-end projects across different
  technologies, with a strong interest in automation, AI-driven workflows, and
  scalable system design. Currently expanding his understanding of enterprise
  technology and systems that make processes more efficient and reliable.
\end{minipage}

\medskip\noindent\color{rule!25}\hrule\color{black}\medskip

\noindent
\begin{minipage}[t]{0.30\linewidth}
  \vspace{0pt}
  \imgOrPlaceholder[width=\linewidth]{../D3/assets/ali_photo.png}{Photo}
\end{minipage}%
\hspace{1em}%
\begin{minipage}[t]{0.66\linewidth}
  \vspace{0pt}
  {\bfseries\color{primary} Ali Aan Khowaja}\\[-2pt]
  {\small\color{muted} Lead Developer / UI Designer}\\[4pt]
  An AI Engineer driven by a simple goal: building applications that work at
  scale. His expertise lies at the intersection of high-performance backend
  architecture and agentic systems. He specialises in designing and deploying
  scalable infrastructures that power intelligent applications. Beyond the code,
  he is a robotics enthusiast passionate about bringing new ideas to life.
\end{minipage}

\subsection{Team Agreement}

The team agreement below was established in Deliverable~1 and remained in force
unchanged throughout the project.

\subsubsection*{Communication}
\begin{itemize}
  \item \textbf{WhatsApp} --- primary channel for daily coordination; respond within 1 hour.
  \item \textbf{In-person} --- preferred when on campus.
  \item \textbf{Google Meet} --- for remote meetings.
  \item Urgent issues: tag the team and respond within 20 minutes.
  \item If unavailable: post an async update in WhatsApp within 24 hours.
\end{itemize}

\subsubsection*{Meetings}
\begin{itemize}
  \item Mandatory: Sprint Planning, Sprint Review, Sprint Retrospective.
  \item Regular check-ins: 4$\times$ per week (in-person or Google Meet).
  \item Absence: notify team 24 hours in advance and share an async update.
  \item Every meeting ends with: decisions + action items + owners.
  \item Notes: rotating assigned person.
\end{itemize}

\subsubsection*{Meeting Preparation}
\begin{itemize}
  \item \textbf{Sprint Planning:} review backlog, propose task breakdowns.
  \item \textbf{Review:} prepare demo flow + screenshots.
  \item \textbf{Retrospective:} each member prepares (1)~what went well,
        (2)~what did not, (3)~what to change.
\end{itemize}

\subsubsection*{Version Control}
\begin{itemize}
  \item Small, atomic commits; new branch per ticket.
  \item PR required before merging; do not commit secrets, credentials, or large binaries.
  \item Commit message format: \texttt{feat:}, \texttt{fix:}, \texttt{docs:}, \texttt{chore:}
\end{itemize}

\subsubsection*{Division of Work (Contract)}
\begin{center}
\begin{tabular}{@{}p{0.20\linewidth}p{0.74\linewidth}@{}}
  \toprule
  \textbf{Owner} & \textbf{Responsibility} \\
  \midrule
  Wajih     & Requirements ownership: user stories, structured specs, and consistency checks. Also: testing and QA. \\[3pt]
  Ali       & Lead UI implementation and feature delivery. \\[3pt]
  Ibraheem  & Coordination, sprint tracking, packaging, and final submission. \\
  \bottomrule
\end{tabular}
\end{center}

\subsubsection*{Contingency Plan}
If a member becomes unavailable or falls behind, the team will re-scope to the
minimum required deliverable content, redistribute work immediately, and
prioritise required artefacts. The PM/SM updates the board, assigns new owners,
and schedules an urgent meeting to unblock progress.

\subsubsection*{Submission Protocol}
\begin{itemize}
  \item \textbf{Submission owner:} Wajih.
  \item \textbf{Internal review:} Ibraheem and Ali.
  \item \textbf{Internal cutoff:} finalise 24~hours before the deadline.
\end{itemize}

% ═══════════════════════════════════════════════════════════════
%  SECTION 4 — INTENDED USE OF THE SYSTEM
% ═══════════════════════════════════════════════════════════════
\section{Intended Use of the System}

\subsection{Stakeholders}

\subsubsection*{Students}
\begin{itemize}
  \item Register and log in with a verified account; choose Student role at registration.
  \item Enrol in a course using a unique class code shared by the instructor.
  \item View enrolled courses, announcements, and deadlines on a course home page.
  \item Submit assignments via file upload; see submission status (Submitted / Late / Missing / Not Submitted).
  \item View grades and instructor feedback in a personal gradebook.
  \item Track all pending work across courses on the To-Do list, with urgency colour-coding.
  \item View per-course assignment status (Assignments tab) and the course roster (People tab).
\end{itemize}

\subsubsection*{Instructors}
\begin{itemize}
  \item Register and log in with a verified account; choose Instructor role at registration.
  \item Create and manage courses; share a system-generated unique class code with students.
  \item Post announcements with optional file and link attachments.
  \item Create assignments with title, description, total points, due date/time, and configurable late policy.
  \item Grade submissions with numeric scores and written feedback.
  \item View a class-wide gradebook (students $\times$ assignments grid).
  \item View the course roster (People tab) with role badges.
\end{itemize}

\subsubsection*{System (Cross-Cutting)}
\begin{itemize}
  \item Enforces role-based access control at every protected action.
  \item Logs important events (logins, submissions, grading updates, enrolment changes) for auditability.
  \item Computes submission status from server-side timestamps, not client time.
\end{itemize}

\subsection{How the System Helps Users}
\begin{itemize}
  \item Consolidates course operations (content, announcements, assessments, grading) into one unified workflow.
  \item Reduces instructor workload by centralising submission tracking and grading.
  \item Provides traceability: all academic actions are recorded and auditable (who submitted what and when; who graded what and when).
  \item Gives students proactive deadline awareness, so missing work is surfaced rather than buried.
\end{itemize}

% ═══════════════════════════════════════════════════════════════
%  SECTION 5 — FEATURES AND OVERALL FUNCTIONALITY
% ═══════════════════════════════════════════════════════════════
\section{Features and Overall Functionality}

{\sloppy
The full feature set is grouped below by module. Each module is mapped to its
backing user stories from Section~\ref{sec:stories}.

\subsection{Authentication and Role-Based Access Control}
\textit{Backed by US\,1--4.}
\begin{itemize}
  \item User self-registration with mandatory role selection (Student or Instructor).
  \item Login with email and password; SHA-256 password hashing at registration and verification.
  \item Logout that explicitly clears the session singleton and returns to the login screen.
  \item RBAC enforced across every protected controller: students never reach
        instructor-only flows (Create Assignment, Post Announcement, Grade
        Submissions). Instructor-only buttons are conditionally rendered ---
        not merely disabled --- so they are invisible to students.
\end{itemize}

\subsection{Course Management and Enrolment}
\textit{Backed by US\,5--9.}
\begin{itemize}
  \item Instructors create courses with a name and optional description.
  \item System generates a unique alphanumeric class code per course at creation.
  \item Students enrol by entering the class code from the catalog view.
  \item Students see only courses they are enrolled in; instructors see only
        courses they own (\texttt{courses.created\_by\,=\,userId}).
  \item Course home page acts as a stream of announcements and assignments.
\end{itemize}

\subsection{Announcements and Attachments}
\textit{Backed by US\,10--13.}
\begin{itemize}
  \item Instructors post announcements with title, body text, and basic formatting.
  \item Optional file attachments via the system FileChooser.
  \item Optional URL attachments for external resources.
  \item Both attachment types persist in the \texttt{announcement\_attachments} table.
  \item Course stream shows attachment chips inline: file chips open via
        \texttt{Desktop.getDesktop()}; link chips open in the default browser.
  \item A single announcement can carry one file and one link attachment independently.
\end{itemize}

\subsection{Assignments and Submissions}
\textit{Backed by US\,14--18.}
\begin{itemize}
  \item Instructors create assignments with title, description, total points,
        due date (DatePicker), due time (hour/minute spinners), and late
        policy toggle (\emph{Allow late turn-ins} vs.\ \emph{Hard deadline}).
  \item Students submit work via the FileChooser within the assignment view.
  \item Server-side timestamp recorded on every submission event.
  \item Automatic status derivation: \emph{Submitted} (before deadline),
        \emph{Late} (after deadline if policy allows), \emph{Missing}
        (deadline passed without upload), \emph{Not Submitted} (deadline
        still in future).
  \item Conditional unsubmit/resubmit based on deadline and instructor policy.
  \item Threaded comment area within the assignment detail view for both students and instructors.
\end{itemize}

\subsection{Grading and Gradebook}
\textit{Backed by US\,19--22.}
\begin{itemize}
  \item Two-panel \texttt{GradingView} per assignment: ListView of submissions
        on the left; numeric grade and feedback editor on the right.
  \item Numeric grade validation against \texttt{total\_points} before save;
        non-numeric input or grade outside $[0, \texttt{total\_points}]$ is rejected inline.
  \item Save issues a parameterised UPDATE on \texttt{submissions.grade} and
        \texttt{submissions.feedback}.
  \item Student personal gradebook lists every graded item with score,
        maximum, and feedback.
  \item Instructor class-wide gradebook is a students $\times$ assignments
        grid backed by a single-JOIN query (\texttt{fetchGradebookForCourse}),
        avoiding N+1 round-trips.
\end{itemize}

\subsection{Course Home: Stream, Assignments, People}
\textit{Backed by US\,9 (Stream), US\,25 (Assignments tab), US\,26 (People tab).}
\begin{itemize}
  \item \textbf{Stream:} reverse-chronological feed of announcements and assignments.
  \item \textbf{Assignments tab:} all course assignments with role-adaptive
        cards. Students see status badges (Submitted / Late / Missing / Not
        Submitted); instructors see a Grade Submissions button per card.
        Status is derived via a LEFT JOIN in
        \texttt{Assignment.fetchWithStatusForUser()}.
  \item \textbf{People tab:} course roster grouped by role badge (Instructor /
        Student), backed by \texttt{Course.fetchEnrolledUsers()}.
\end{itemize}

\subsection{To-Do List and Deadline Awareness}
\textit{Backed by US\,23--24.}
\begin{itemize}
  \item Aggregates pending assignments across every enrolled course.
  \item Partitioned into \emph{Active} (deadline in future) and \emph{Missing} (deadline passed without submission).
  \item Urgency colour-coding: red (overdue), orange (due within 48~hours), default (more than 48~hours).
  \item Click an item to open the assignment detail view.
  \item Items are removed from the active list automatically once submitted.
\end{itemize}
\par}

% ═══════════════════════════════════════════════════════════════
%  SECTION 6 — USER STORIES
% ═══════════════════════════════════════════════════════════════
\section{Master User Stories}
\label{sec:stories}

The system was scoped against 26 user stories: 24 from the original product
backlog established in Deliverable~1, plus US\,25 and US\,26 (Assignments tab
and People tab) which were identified during Sprint~3 planning when the
course-home navigation gap became apparent.

\subsection{Authentication and Role-Based Access Control}
\begin{enumerate}[label=\textbf{US\,\arabic*}]
  \item As a \emph{User} (Student/Instructor), I want to log in securely so
        that I can access my account and authorised system features.
  \item As a \emph{User}, I want to log out of my account so that my session
        is secure when I am not using the system.
  \item As a \emph{User}, I want to select my role (Student or Instructor)
        during registration so that I can use the platform appropriately.
  \item As a \emph{User}, I want role-based access control to be enforced
        across all features so that individuals can only perform actions
        permitted by their chosen role.
\end{enumerate}

\subsection{Course Management and Enrolment}
\begin{enumerate}[label=\textbf{US\,\arabic*}, start=5]
  \item As an \emph{Instructor}, I want to create a new course so that I have
        a dedicated space to manage learning materials and track student
        progress.
  \item As a \emph{User}, I want to view a course catalogue/listing so that I
        can see the courses I am enrolled in or teaching.
  \item As an \emph{Instructor}, I want the system to generate a unique random
        code for my course so that I can share it securely with my students.
  \item As a \emph{Student}, I want to enrol in a course using a unique class
        code so that I can self-register for the classes I am supposed to take.
  \item As a \emph{Student}, I want to view a course home page that displays
        announcements and upcoming items so that I am immediately aware of
        current course activities.
\end{enumerate}

\subsection{Announcements and Course Material}
\begin{enumerate}[label=\textbf{US\,\arabic*}, start=10]
  \item As an \emph{Instructor}, I want to post announcements so that I can
        broadcast important updates or instructions to all enrolled students.
  \item As a \emph{Student}, I want to see an announcement feed for my course
        so that I never miss critical information from my instructor.
  \item As an \emph{Instructor}, I want to attach a file to an announcement so
        that I can share reference materials directly from the course stream.
  \item As an \emph{Instructor}, I want to attach a URL to an announcement so
        that I can share links to external resources without requiring a file
        upload.
\end{enumerate}

\subsection{Assignments and Submissions}
\begin{enumerate}[label=\textbf{US\,\arabic*}, start=14]
  \item As an \emph{Instructor}, I want to create assignments with a title,
        description, and due date/time so that expectations are clear to
        students.
  \item As an \emph{Instructor}, I want to define late policies and
        resubmission rules so that I can enforce grading standards.
  \item As a \emph{Student}, I want to submit my assignments via file upload
        so that I can turn in my completed work electronically.
  \item As an \emph{Instructor}, I want the system to timestamp all
        submissions so that I can accurately determine their status (submitted,
        late, missing).
  \item As a \emph{Student}, I want to see the status of my submissions (e.g.,
        submitted on time, late, missing) so that I know my work has been
        received.
\end{enumerate}

\subsection{Grading, Feedback, and Gradebook}
\begin{enumerate}[label=\textbf{US\,\arabic*}, start=19]
  \item As an \emph{Instructor}, I want to view a list of student submissions
        for an assignment so that I can review them efficiently.
  \item As an \emph{Instructor}, I want to assign scores and add comments to
        student submissions so that I can provide them with constructive
        feedback.
  \item As an \emph{Instructor}, I want to view a comprehensive gradebook of
        all students and their graded items so that I can monitor the overall
        performance of the class.
  \item As a \emph{Student}, I want to view my personal gradebook with scores
        and instructor feedback so that I can track my academic standing and
        progress.
\end{enumerate}

\subsection{To-Do List and Deadlines}
\begin{enumerate}[label=\textbf{US\,\arabic*}, start=23]
  \item As a \emph{Student}, I want to view a to-do list within my course
        catalogue section so that I have complete visibility over my progress
        (submitted vs.\ missing items).
  \item As a \emph{Student}, I want the system to highlight upcoming due dates
        and late/missed items so that I can effectively manage my deadlines and
        prioritise my workload.
\end{enumerate}

\subsection{Course Home Navigation \texorpdfstring{\normalfont\small(Added in Sprint 3)}{}}
\begin{enumerate}[label=\textbf{US\,\arabic*}, start=25]
  \item As a \emph{User}, I want a dedicated Assignments tab in the course
        home showing all assignments and their status at a glance, so I do
        not have to navigate away from the course to see my workload.
  \item As a \emph{User}, I want a People tab in the course home showing
        everyone enrolled in the course with their role badge, so I know who
        is in my class.
\end{enumerate}

% ═══════════════════════════════════════════════════════════════
%  SECTION 7 — STRUCTURED SPECIFICATIONS
% ═══════════════════════════════════════════════════════════════
\section{Structured Specifications}
\label{sec:specs}

The structured specifications below define each feature module's behaviour at
a granularity that makes acceptance testable. Each spec is grouped by the
module it covers in Section~\ref{sec:stories}.

% ── Spec: Auth
\begin{userstory}
  \textbf{Spec 1 --- US\,1--4: Authentication, Registration, and RBAC}
\end{userstory}
\begin{center}
\begin{tabular}{@{}p{0.22\linewidth}@{\hspace{6pt}}p{0.72\linewidth}@{}}
  \toprule
  \fcol{Field} & Detail \\
  \midrule
  \fcol{Actor} & Any user (Student or Instructor); System (for RBAC enforcement) \\
  \fcol{Inputs} & \emph{Register:} username, email, password, role (Student / Instructor). \emph{Login:} email, password. \\
  \fcol{Processing} & Registration: validate all fields non-empty; hash password with SHA-256; \texttt{User.register()} inserts the user and returns an empty \texttt{Optional} on duplicate username/email. Login: hash submitted password and compare to stored hash; populate \texttt{Session} on match. Logout: \texttt{Session.logout()} clears state; navigate to \texttt{LoginView}. RBAC: every protected controller checks \texttt{Session.getInstance().getRole()} before executing instructor-only logic. \\
  \fcol{Outputs} & On register/login success: user lands on \texttt{CourseCatalogView} with an active session. On logout: session cleared; user returns to \texttt{LoginView}. RBAC: students never see instructor-only buttons (conditionally rendered, not disabled). \\
  \fcol{Preconditions} & Register: username and email not already in \texttt{users}. Login: account exists with a stored hash. RBAC: user is logged in. \\
  \fcol{Postconditions} & \texttt{Session} holds \texttt{id, username, email, role}. Logout leaves \texttt{Session} empty; no credentials persist. No student can trigger an instructor action even by direct controller call. \\
  \fcol{Exceptions} & Wrong password / unknown email $\rightarrow$ ``Incorrect email or password''. Empty fields $\rightarrow$ ``Please fill in all fields''. Duplicate username/email $\rightarrow$ ``Username or email is already taken''. \\
  \bottomrule
\end{tabular}
\end{center}

% ── Spec: Course Mgmt
\begin{userstory}
  \textbf{Spec 2 --- US\,5--9: Course Creation, Catalogue, and Class-Code Enrolment}
\end{userstory}
\begin{center}
\begin{tabular}{@{}p{0.22\linewidth}@{\hspace{6pt}}p{0.72\linewidth}@{}}
  \toprule
  \fcol{Field} & Detail \\
  \midrule
  \fcol{Actor} & Instructor (create); Student (enrol, view); Both (catalogue) \\
  \fcol{Inputs} & \emph{Create course:} name (required), description (optional). \emph{Enrol:} class code. \\
  \fcol{Processing} & On course creation: generate a unique alphanumeric class code; persist with \texttt{created\_by\,=\,Session.userId}. On enrol: look up course by code; insert into \texttt{course\_enrollments}. Catalogue: students see courses joined to \texttt{course\_enrollments}; instructors see courses where \texttt{created\_by\,=\,userId}. \\
  \fcol{Outputs} & New course visible on instructor's catalogue. Enrolled course visible on student's catalogue. Course home displays a stream of recent announcements, latest assignments, and immediate deadlines. \\
  \fcol{Preconditions} & User is logged in. \\
  \fcol{Postconditions} & Course row persisted with unique code. Enrolment row persisted linking student to course. \\
  \fcol{Exceptions} & Invalid or unknown class code $\rightarrow$ inline error ``Course not found''. Already enrolled $\rightarrow$ ``Already enrolled''. \\
  \bottomrule
\end{tabular}
\end{center}

% ── Spec: Announcements + Attachments
\begin{userstory}
  \textbf{Spec 3 --- US\,10--13: Announcements and Attachments}
\end{userstory}
\begin{center}
\begin{tabular}{@{}p{0.22\linewidth}@{\hspace{6pt}}p{0.72\linewidth}@{}}
  \toprule
  \fcol{Field} & Detail \\
  \midrule
  \fcol{Actor} & Instructor (post); all enrolled users (view) \\
  \fcol{Inputs} & Title (required), body (required), optional file via FileChooser, optional URL. \\
  \fcol{Processing} & Validate title and body non-empty; insert announcement with author ID, course ID, and server timestamp. Attachments: file row (\texttt{is\_link\,=\,false}, original filename + path) or link row (\texttt{is\_link\,=\,true}, \texttt{link\_url}) in \texttt{announcement\_attachments}, foreign-keyed to the announcement. A single announcement may carry one file and one link independently. \\
  \fcol{Outputs} & Announcement appears in the course stream (newest first) for all enrolled students. Attachment chips render inline beneath the body: file chip opens via \texttt{Desktop.getDesktop()}; link chip opens in the default browser. \\
  \fcol{Preconditions} & Instructor is logged in and owns the course. \\
  \fcol{Postconditions} & Announcement row persisted; attachment rows (if any) persisted; visible on next stream load. \\
  \fcol{Exceptions} & Empty title or body $\rightarrow$ inline validation error; form does not submit. No file/URL provided $\rightarrow$ announcement saves with no attachment row. \\
  \bottomrule
\end{tabular}
\end{center}

% ── Spec: Assignments + Submissions
\begin{userstory}
  \textbf{Spec 4 --- US\,14--18: Assignment Creation, Submission, and Status Tracking}
\end{userstory}
\begin{center}
\begin{tabular}{@{}p{0.22\linewidth}@{\hspace{6pt}}p{0.72\linewidth}@{}}
  \toprule
  \fcol{Field} & Detail \\
  \midrule
  \fcol{Actor} & Instructor (create); Student (submit); System (status derivation) \\
  \fcol{Inputs} & \emph{Create:} title, description (optional), total points (numeric), due date, due time, late policy toggle. \emph{Submit:} file via FileChooser. \\
  \fcol{Processing} & On create: validate title, points, deadline; combine date and time into a \texttt{LocalDateTime}; persist linked to course with late policy flag. On submit: record server-side timestamp; compare to deadline; evaluate late policy; tag status. Status derivation in \texttt{Assignment.fetchWithStatusForUser()} via LEFT JOIN: \emph{Submitted} (before deadline), \emph{Late} (after deadline, policy permits), \emph{Missing} (deadline passed, no upload), \emph{Not Submitted} (deadline still in future). \\
  \fcol{Outputs} & Assignment appears in course stream and on enrolled students' To-Do lists. Submission row created/updated with file path and timestamp; status badge visible to student and instructor. \\
  \fcol{Preconditions} & Create: instructor logged in and owns course. Submit: student is enrolled; assignment has a deadline. \\
  \fcol{Postconditions} & Assignment persisted with policy flag. Submission row stored or updated; status visible. \\
  \fcol{Exceptions} & Missing title / deadline / non-numeric points $\rightarrow$ validation error. No file selected $\rightarrow$ Submit button inactive. Late submission with hard policy $\rightarrow$ submission blocked. \\
  \bottomrule
\end{tabular}
\end{center}

% ── Spec: Grading + Gradebook
\begin{userstory}
  \textbf{Spec 5 --- US\,19--22: Grading Interface and Gradebook}
\end{userstory}
\begin{center}
\begin{tabular}{@{}p{0.22\linewidth}@{\hspace{6pt}}p{0.72\linewidth}@{}}
  \toprule
  \fcol{Field} & Detail \\
  \midrule
  \fcol{Actor} & Instructor (grade, class view); Student (personal view) \\
  \fcol{Inputs} & Grading: selected submission, numeric grade (required), feedback text (optional). Personal gradebook: none beyond authenticated session. \\
  \fcol{Processing} & Grading: load all submissions via \texttt{Submission.fetchByAssignmentId()}; pre-fill if previously graded; validate grade is numeric and within $0$--\texttt{total\_points}; \texttt{Submission.saveGrade()} issues a parameterised UPDATE. Gradebook: \texttt{fetchGradedForStudent} for personal view; single-JOIN \texttt{fetchGradebookForCourse} for instructor class grid. \\
  \fcol{Outputs} & \texttt{submissions.grade} and \texttt{submissions.feedback} updated; ListView reflects graded state; student gradebook shows score and feedback immediately. Instructor sees students $\times$ assignments grid. \\
  \fcol{Preconditions} & Instructor is logged in and owns the course; assignment exists. \\
  \fcol{Postconditions} & Submission row has non-null \texttt{grade} and optionally non-null \texttt{feedback}. Original submission timestamp preserved. \\
  \fcol{Exceptions} & No submissions $\rightarrow$ empty ListView with empty-state label. Non-numeric grade or grade $>$ \texttt{total\_points} $\rightarrow$ inline validation error; Save blocked. \\
  \bottomrule
\end{tabular}
\end{center}

% ── Spec: To-Do
\begin{userstory}
  \textbf{Spec 6 --- US\,23--24: To-Do List and Deadline Awareness}
\end{userstory}
\begin{center}
\begin{tabular}{@{}p{0.22\linewidth}@{\hspace{6pt}}p{0.72\linewidth}@{}}
  \toprule
  \fcol{Field} & Detail \\
  \midrule
  \fcol{Actor} & Student \\
  \fcol{Inputs} & System reads all assignments across enrolled courses on view load. \\
  \fcol{Processing} & Fetch all assignments for enrolled courses; filter out submitted items; classify remaining as \emph{Active} (deadline in future) or \emph{Missing} (deadline passed without submission); compute urgency relative to current server time. \\
  \fcol{Outputs} & To-Do view with Active and Missing sections; items colour-coded red (overdue), orange (due within 48~hours), default (more than 48~hours away); clicking opens AssignmentDetails. \\
  \fcol{Preconditions} & Student is logged in and enrolled in at least one course. \\
  \fcol{Postconditions} & List reflects current state. Items disappear from Active automatically on submission. \\
  \fcol{Exceptions} & No pending assignments $\rightarrow$ empty-state message. \\
  \bottomrule
\end{tabular}
\end{center}

% ═══════════════════════════════════════════════════════════════
%  NON-FUNCTIONAL REQUIREMENTS
% ═══════════════════════════════════════════════════════════════
\section{Non-Functional Requirements}

\begin{center}
\begin{tabular}{@{}p{0.22\linewidth}@{\hspace{6pt}}p{0.72\linewidth}@{}}
  \toprule
  \fcol{NFR Category} & Requirement \\
  \midrule
  Security
    & Passwords are never stored or compared in plain text --- SHA-256
      hashing via \texttt{PasswordUtil.hash()} at all registration and
      login points. Session state held in a singleton, explicitly cleared on
      logout. RBAC enforced at the controller level before every
      instructor-only action. All SQL is parameterised; SQL injection at the
      login boundary is blocked (verified by TC-06). User inputs validated
      and sanitised. \\[6pt]
  Performance
    & Primary pages (course home, assignment list, gradebook, To-Do list)
      load within 2--3~seconds at typical class sizes ($\sim$200 students per
      course) on lab hardware. Class-wide gradebook uses a single JOIN, not
      per-student queries. File uploads complete within reasonable time for
      typical submission sizes (documents and PDFs up to 10~MB). \\[6pt]
  Reliability
    & Submissions are not lost once acknowledged. Grade saves are atomic
      (single parameterised UPDATE; no partial state). Submission status is
      computed from server-side timestamps, not client time. Session logout
      completes synchronously before the login screen is shown. \\[6pt]
  Usability
    & Consistent navigation and layout across all views. Status badges
      (Submitted / Late / Missing / Not Submitted) use both colour and text
      so they are not colour-only. To-Do list urgency colours are consistent
      with the status badge palette. Gradebook cells show ``---'' for
      ungraded items rather than blank or null. Instructor-only buttons are
      conditionally rendered for students, not merely disabled. \\[6pt]
  Maintainability
    & Layered MVC architecture (controller / model / view) enforced across
      all features. Business logic (deadline evaluation, status tagging,
      grade validation, RBAC checks) lives in model and controller classes,
      not in FXML event handlers. New views follow the same FXML + CSS
      structure as existing ones. \\[6pt]
  Scalability
    & Supports multiple courses and concurrent users typical of a
      department-scale deployment (not enterprise scale). The JOIN-based
      gradebook query and LEFT JOIN status derivation avoid per-student
      round-trips that would not scale past a single classroom. \\[6pt]
  Auditability
    & Important events (logins, submissions, grading updates, enrolment
      changes) are logged with server timestamps for traceability. Submission
      events (upload, unsubmit, resubmit) are logged. Original submission
      timestamps are preserved alongside grade and feedback fields --- no
      information is overwritten on grading. \\
  \bottomrule
\end{tabular}
\end{center}

% ═══════════════════════════════════════════════════════════════
%  SECTION 10 — TESTING & QA
% ═══════════════════════════════════════════════════════════════
\section{Testing and Quality Assurance}
\label{sec:testing}

\begingroup
\setlength{\tabcolsep}{4pt}

\subsection{Software Test Plan}

\begin{center}
\begin{tabular}{@{}p{0.22\linewidth}@{\hspace{6pt}}p{0.72\linewidth}@{}}
  \toprule
  \fcol{Field} & Detail \\
  \midrule
  \fcol{Scope} & Authentication (register/login/logout), RBAC, Announcement Attachments, Assignments Tab, People Tab, Grading Interface, Gradebook, plus regression checks on Sprint~1 and~2 features. \\
  \fcol{Objectives} & Verify each feature meets its acceptance criteria; detect regressions; validate RBAC enforcement at every protected action. \\
  \fcol{Test Types} & Black-box functional testing (equivalence partitioning, boundary value analysis, error guessing); white-box structural testing (branch coverage on critical controllers). \\
  \fcol{Test Environment} & Windows~11, Java~24, JavaFX~21, PostgreSQL~16; seeded via \texttt{data/seed.sql} (3 users: \texttt{ali} = instructor, \texttt{ibbi} + \texttt{sara} = students; 2 courses; 6 assignments; 6 submissions in mixed states). \\
  \fcol{Entry Criteria} & \texttt{make seed} succeeds; \texttt{make compile} produces zero errors; application launches to \texttt{LoginView}. \\
  \fcol{Exit Criteria} & All P1 (critical) test cases pass; no unfixed P1 bugs; all identified P2 bugs documented with status. \\
  \fcol{Tester} & Wajih-Ur-Raza Asif. \\
  \bottomrule
\end{tabular}
\end{center}

\subsection{Black-Box Test Cases}

\subsubsection*{Module: Authentication --- Login}
\begin{center}\small
\begin{tabular}{@{}p{0.06\linewidth}p{0.13\linewidth}p{0.30\linewidth}p{0.33\linewidth}p{0.08\linewidth}@{}}
  \toprule
  \textbf{TC \#} & \textbf{Technique} & \textbf{Input} & \textbf{Expected Output} & \textbf{Status} \\
  \midrule
  TC-01 & Valid partition   & \texttt{ali@naawbi.edu} / \texttt{password123}    & Login succeeds; catalog loads; role = instructor  & \textcolor{success}{Pass} \\
  TC-02 & Valid partition   & \texttt{ibbi@naawbi.edu} / \texttt{password123}   & Login succeeds; catalog loads; role = student     & \textcolor{success}{Pass} \\
  TC-03 & Invalid partition & \texttt{ali@naawbi.edu} / \texttt{wrongpass}      & Inline error ``Incorrect email or password''      & \textcolor{success}{Pass} \\
  TC-04 & Invalid partition & \texttt{nobody@naawbi.edu} / \texttt{password123} & Inline error ``Incorrect email or password''      & \textcolor{success}{Pass} \\
  TC-05 & Boundary / empty  & Empty email + empty password                      & Inline error ``Please fill in all fields''        & \textcolor{success}{Pass} \\
  TC-06 & Error guessing    & \texttt{' OR 1=1 -{}-} / any                      & Login fails --- parameterised query blocks SQLi   & \textcolor{success}{Pass} \\
  \bottomrule
\end{tabular}
\end{center}

\subsubsection*{Module: Authentication --- Registration}
\begin{center}\small
\begin{tabular}{@{}p{0.06\linewidth}p{0.13\linewidth}p{0.30\linewidth}p{0.33\linewidth}p{0.08\linewidth}@{}}
  \toprule
  \textbf{TC \#} & \textbf{Technique} & \textbf{Input} & \textbf{Expected Output} & \textbf{Status} \\
  \midrule
  TC-07 & Valid partition   & New unique username, email, password, role=Student & Account created; redirected to catalog       & \textcolor{success}{Pass} \\
  TC-08 & Invalid partition & Username already taken                            & ``Username or email is already taken''        & \textcolor{success}{Pass} \\
  TC-09 & Invalid partition & Email already taken                                & ``Username or email is already taken''        & \textcolor{success}{Pass} \\
  TC-10 & Boundary / empty  & Any field left blank                               & ``Please fill in all fields''                 & \textcolor{success}{Pass} \\
  \bottomrule
\end{tabular}
\end{center}

\subsubsection*{Module: RBAC --- Instructor-Only Guard}
\begin{center}\small
\begin{tabular}{@{}p{0.06\linewidth}p{0.13\linewidth}p{0.30\linewidth}p{0.33\linewidth}p{0.08\linewidth}@{}}
  \toprule
  \textbf{TC \#} & \textbf{Technique} & \textbf{Input} & \textbf{Expected Output} & \textbf{Status} \\
  \midrule
  TC-11 & Invalid partition & Student \texttt{ibbi} clicks ``Create Assignment'' & Access Denied alert; no form opens          & \textcolor{success}{Pass} \\
  TC-12 & Valid partition   & Instructor \texttt{ali} clicks ``Create Assignment'' & Create Assignment form opens normally       & \textcolor{success}{Pass} \\
  TC-13 & Valid partition   & Student logs in                                  & Catalog shows only enrolled courses (CS101, CS202) & \textcolor{success}{Pass} \\
  TC-14 & Valid partition   & Instructor logs in                               & Catalog shows only courses created by \texttt{ali}   & \textcolor{success}{Pass} \\
  \bottomrule
\end{tabular}
\end{center}

\subsubsection*{Module: Grading Interface}
\begin{center}\small
\begin{tabular}{@{}p{0.06\linewidth}p{0.13\linewidth}p{0.30\linewidth}p{0.33\linewidth}p{0.08\linewidth}@{}}
  \toprule
  \textbf{TC \#} & \textbf{Technique} & \textbf{Input} & \textbf{Expected Output} & \textbf{Status} \\
  \midrule
  TC-15 & Valid partition & grade=85 / total=100 / feedback=``Good work'' & Grade saved; submission shows graded in ListView & \textcolor{success}{Pass} \\
  TC-16 & Boundary value  & grade = 0                                     & Grade saved as 0; valid                           & \textcolor{success}{Pass} \\
  TC-17 & Boundary value  & grade = total\_points (e.g.\ 100)             & Grade saved at maximum; valid                     & \textcolor{success}{Pass} \\
  TC-18 & Boundary value  & grade = total\_points + 1 (e.g.\ 101)         & Inline validation error; Save blocked             & \textcolor{success}{Pass} \\
  TC-19 & Invalid partition & grade = ``abc'' (non-numeric)               & Inline validation error; Save blocked             & \textcolor{success}{Pass} \\
  TC-20 & Valid partition & No feedback; grade = 50                       & Grade saved with NULL feedback                    & \textcolor{success}{Pass} \\
  \bottomrule
\end{tabular}
\end{center}

\subsubsection*{Module: Submission Status Derivation}
\begin{center}\small
\begin{tabular}{@{}p{0.06\linewidth}p{0.13\linewidth}p{0.30\linewidth}p{0.33\linewidth}p{0.08\linewidth}@{}}
  \toprule
  \textbf{TC \#} & \textbf{Technique} & \textbf{Input} & \textbf{Expected Output} & \textbf{Status} \\
  \midrule
  TC-21 & Valid partition   & Submission before deadline                        & Badge: \textbf{Submitted}       & \textcolor{success}{Pass} \\
  TC-22 & Valid partition   & Submission after deadline; late policy allowed   & Badge: \textbf{Late}            & \textcolor{success}{Pass} \\
  TC-23 & Valid partition   & Deadline passed; no submission; late allowed      & Badge: \textbf{Missing}         & \textcolor{success}{Pass} \\
  TC-24 & Valid partition   & Deadline in future; no submission                 & Badge: \textbf{Not Submitted}   & \textcolor{success}{Pass} \\
  TC-25 & Invalid partition & Deadline passed; hard policy; attempt submit      & Submission blocked              & \textcolor{success}{Pass} \\
  \bottomrule
\end{tabular}
\end{center}

\subsection{White-Box Testing}

\subsubsection*{LoginController.handleLogin() --- Branch Coverage}
\begin{quote}\small\ttfamily
B1: email or password empty \hfill $\to$ show error, return\\
B2: User.findByEmail() returns empty \hfill $\to$ show error, return\\
B3: stored hash does not match input hash \hfill $\to$ show error, return\\
B4: all checks pass \hfill $\to$ populate Session, navigate
\end{quote}

\begin{center}
\begin{tabular}{@{}p{0.28\linewidth}p{0.56\linewidth}p{0.10\linewidth}@{}}
  \toprule
  \fcolplain{Branch} & Test Case & Covered \\
  \midrule
  \fcolplain{B1 --- empty fields} & TC-05                  & Yes \\
  \fcolplain{B2 --- unknown email} & TC-04                  & Yes \\
  \fcolplain{B3 --- wrong password} & TC-03                  & Yes \\
  \fcolplain{B4 --- successful login} & TC-01, TC-02           & Yes \\
  \midrule
  \multicolumn{2}{r}{\textbf{Branch coverage}} & \textbf{4\,/\,4 (100\%)} \\
  \bottomrule
\end{tabular}
\end{center}

\subsubsection*{GradingController.handleSave() --- Branch Coverage}
\begin{quote}\small\ttfamily
B1: grade empty or non-numeric \hfill $\to$ show error, return\\
B2: grade $>$ total\_points \hfill $\to$ show error, return\\
B3: grade $<$ 0 \hfill $\to$ show error, return\\
B4: all checks pass \hfill $\to$ Submission.saveGrade(), refresh
\end{quote}

\begin{center}
\begin{tabular}{@{}p{0.28\linewidth}p{0.56\linewidth}p{0.10\linewidth}@{}}
  \toprule
  \fcolplain{Branch} & Test Case & Covered \\
  \midrule
  \fcolplain{B1 --- non-numeric input} & TC-19                        & Yes \\
  \fcolplain{B2 --- grade over maximum} & TC-18                        & Yes \\
  \fcolplain{B3 --- grade below minimum} & Boundary: grade = $-1$ (manual) & Yes \\
  \fcolplain{B4 --- valid save} & TC-15, TC-16, TC-17          & Yes \\
  \midrule
  \multicolumn{2}{r}{\textbf{Branch coverage}} & \textbf{4\,/\,4 (100\%)} \\
  \bottomrule
\end{tabular}
\end{center}

\subsubsection*{Assignment.fetchWithStatusForUser() --- Statement Coverage}
The method executes a LEFT JOIN and derives submission status via conditional
logic. All four status paths (Submitted, Late, Missing, Not Submitted) are
exercised by TC-21 through TC-24 against the seeded dataset --- achieving
statement coverage of every reachable branch in the status derivation block.

\subsection{Automated JUnit Test Suite}

Alongside the manual cases above, the project ships a JUnit~5 test suite
that runs on \texttt{make test}. It is split into two layers.

\textbf{Unit tests} cover the non-database parts of the model.
\texttt{PasswordUtilTest} pins the SHA-256 hashing contract (determinism,
hex format, a known vector for \texttt{password123}, unicode input,
\texttt{null} handling). \texttt{SessionTest} covers the singleton,
login and logout, role checks, and the \texttt{IllegalStateException}
thrown by every getter when no user is logged in. The remaining classes
do round-trip checks on the constructors, getters, and setters of every
model.

\textbf{DAO integration tests} run the real JDBC code in \texttt{User},
\texttt{Course}, and \texttt{Submission} against a dedicated PostgreSQL
database called \texttt{naawbi\_test}, so the development database is
never written to. Each test starts from a truncated schema, so order
does not matter. \texttt{IntegrationTestBase} also refuses to start
unless the JDBC URL contains the string \texttt{naawbi\_test}, as a
safety net against accidental runs against \texttt{naawbi}.

\begin{center}
\small\begin{tabular}{@{}p{0.30\linewidth}p{0.10\linewidth}p{0.50\linewidth}@{}}
  \toprule
  \fcolplain{Layer} & Tests & Notes \\
  \midrule
  Unit (no database)        & 44 & 9 classes, runs in $\sim$0.6\,s \\
  DAO integration           & 17 & 3 classes, runs in $\sim$5\,s, hits \texttt{naawbi\_test} \\
  \midrule
  \multicolumn{2}{r}{\textbf{Total}} & \textbf{61 tests, all passing} \\
  \bottomrule
\end{tabular}
\end{center}

\noindent
Setup for teammates is a single extra step on top of the existing
prerequisites: drop \texttt{junit-platform-console-standalone-1.11.4.jar}
into \texttt{lib/}. The \texttt{naawbi\_test} database is created
automatically the first time \texttt{make test} runs (handled by
\texttt{data/db-test-init.sh}).

\subsection{Bug Report}

\begin{center}\footnotesize\sloppy
\begin{tabular}{@{}p{0.07\linewidth}p{0.11\linewidth}p{0.60\linewidth}p{0.06\linewidth}p{0.07\linewidth}@{}}
  \toprule
  \textbf{Bug ID} & \textbf{Module} & \textbf{Description and Fix} & \textbf{Sev.} & \textbf{Status} \\
  \midrule
  BUG-01 & Grading     & \texttt{GradingView} threw NPE on an assignment with zero submissions. \emph{Fix:} empty-list guard before ListView population.                                            & P1 & \textcolor{success}{Fixed} \\
  BUG-02 & RBAC        & Instructors saw courses they did not create because \texttt{created\_by} was not being set. \emph{Fix:} \texttt{CreateCourseController} now sets it from \texttt{Session}.     & P1 & \textcolor{success}{Fixed} \\
  BUG-03 & Gradebook   & GridPane column headers misaligned when a student had no submissions for the first assignment. \emph{Fix:} iterate all assignments independent of submission presence.        & P2 & \textcolor{success}{Fixed} \\
  BUG-04 & Attachments & Attachment chip did not render after stream refresh. \emph{Fix:} stream loader now calls \texttt{fetchByAnnouncementId()} per announcement. & P1 & \textcolor{success}{Fixed} \\
  \bottomrule
\end{tabular}
\end{center}

\subsection{Test Execution Results}

Testing was performed manually against the seeded database (\texttt{make seed})
using all three test accounts.

\begin{center}
\small\begin{tabular}{@{}p{0.24\linewidth}p{0.12\linewidth}p{0.58\linewidth}@{}}
  \toprule
  \fcolplain{Account} & Role       & Scenarios Tested \\
  \midrule
  \texttt{ali@naawbi.edu}  & Instructor & Login, RBAC guard, create assignment, post announcement with both attachment types, grade all 6 submissions, view class gradebook, People tab, logout. \\
  \texttt{ibbi@naawbi.edu} & Student    & Login, register, enrolment filtering, submit assignment, view submission status, personal gradebook, To-Do list, logout. \\
  \texttt{sara@naawbi.edu} & Student    & Login, second-student perspective, verify gradebook isolation, People tab roster. \\
  \bottomrule
\end{tabular}
\end{center}

\noindent\textbf{Summary.} All 25 black-box test cases passed. All 4 identified
bugs (BUG-01 through BUG-04) were fixed before submission. No unfixed P1 bugs
remain. One P2 issue (BUG-03) was found and fixed. System behaviour matched
acceptance criteria across all delivered user stories.

\subsection{How Testing Improved System Quality}
The four bugs caught in test would all have been user-facing in submission:
BUG-01 would have crashed the Grading modal for any assignment with zero
submissions (a common early state); BUG-02 broke the core RBAC invariant
entirely; BUG-04 made the attachment feature invisible to students. Explicit
coverage of all four submission-status paths also surfaced a gap in the
``Not Submitted'' derivation that was corrected before integration.

\endgroup

% ═══════════════════════════════════════════════════════════════
%  SECTION 11 — IMPLEMENTATION SCREENSHOTS
% ═══════════════════════════════════════════════════════════════
\section{Implementation Screenshots}

Screenshots below are organised by module, showing the system as delivered.

\begin{figure}[H]
  \centering
  \begin{minipage}{0.48\textwidth}\centering
    \imgOrPlaceholder[width=\textwidth]{../D3/assets/login_view.png}{Login View}
    \subcaption{LoginView --- tabbed Sign In / Register.}
  \end{minipage}\hfill
  \begin{minipage}{0.48\textwidth}\centering
    \imgOrPlaceholder[width=\textwidth]{../D3/assets/course_catalog_role.png}{Course Catalog}
    \subcaption{Course Catalog filtered by role (instructor).}
  \end{minipage}
  \caption{Authentication and role-based catalog filtering.}
\end{figure}

\begin{figure}[H]
  \centering
  \imgOrPlaceholder[width=0.72\textwidth]{../D3/assets/announcement_attachments.png}{Announcement Attachments}
  \caption{Course Stream --- announcement with file and link attachment chips stored in \texttt{announcement\_attachments}.}
\end{figure}

\begin{figure}[H]
  \centering
  \begin{minipage}{0.48\textwidth}\centering
    \imgOrPlaceholder[width=\textwidth]{../D3/assets/assignments_tab_student.png}{Assignments Tab — Student}
    \subcaption{Student view with status badges.}
  \end{minipage}\hfill
  \begin{minipage}{0.48\textwidth}\centering
    \imgOrPlaceholder[width=\textwidth]{../D3/assets/assignments_tab_instructor.png}{Assignments Tab — Instructor}
    \subcaption{Instructor view with Grade Submissions buttons.}
  \end{minipage}
  \caption{Assignments Tab --- role-adaptive cards sharing the same FXML backbone.}
\end{figure}

\begin{figure}[H]
  \centering
  \imgOrPlaceholder[width=0.72\textwidth]{../D3/assets/grading_modal.png}{Grading Modal}
  \caption{GradingView --- ListView of submissions on the left; numeric grade and feedback editor on the right.}
\end{figure}

\begin{figure}[H]
  \centering
  \begin{minipage}{0.48\textwidth}\centering
    \imgOrPlaceholder[width=\textwidth]{../D3/assets/student_gradebook.png}{Student Gradebook}
    \subcaption{Student personal gradebook (read-only).}
  \end{minipage}\hfill
  \begin{minipage}{0.48\textwidth}\centering
    \imgOrPlaceholder[width=\textwidth]{../D3/assets/instructor_gradebook.png}{Instructor Class Gradebook}
    \subcaption{Instructor class-wide grid from \texttt{fetchGradebookForCourse()}.}
  \end{minipage}
  \caption{Gradebook views --- role-adaptive rendering (table vs.\ GridPane).}
\end{figure}

\begin{figure}[H]
  \centering
  \imgOrPlaceholder[width=0.72\textwidth]{../D3/assets/people_tab.png}{People Tab}
  \caption{People Tab --- role-grouped roster cards from \texttt{Course.fetchEnrolledUsers()}.}
\end{figure}

% ═══════════════════════════════════════════════════════════════
%  SECTION 12 — DEVELOPMENT METHODOLOGY (SCRUM)
% ═══════════════════════════════════════════════════════════════
\section{Development Methodology (Scrum)}

The project followed a three-sprint Scrum cadence with mandatory Sprint
Planning, Sprint Review, and Sprint Retrospective ceremonies. Issue tracking,
backlog management, and burn-down charts were maintained on GitHub Projects.

\subsection{Sprint Summary}

\begin{center}\small
\begin{tabular}{@{}p{0.10\linewidth}p{0.20\linewidth}p{0.34\linewidth}p{0.28\linewidth}@{}}
  \toprule
  \fcol{Sprint} & \textbf{Duration} & \textbf{Sprint Goal} & \textbf{Modules Delivered} \\
  \midrule
  Sprint 1 & Feb 2026
           & Establish course structure: course creation, catalogue, and class-code enrolment.
           & Course Creation, Course Catalogue + Home, Class-Code Enrolment. \\[6pt]
  Sprint 2 & 1 -- 25 March 2026
           & Implement course announcements, assignment creation with deadlines and late policies, file-based submission tracking, and the student To-Do list.
           & Announcements, Assignment Creation, File Submissions \& Status, To-Do List. \\[6pt]
  Sprint 3 & 26 March -- 20 April 2026
           & Close the platform's remaining gaps: identity, access, and evaluation.
           & Authentication, RBAC, Announcement Attachments, Assignments Tab, People Tab, Grading, Gradebook. \\
  \bottomrule
\end{tabular}
\end{center}

\noindent As a planning note worth recording: authentication was originally
scoped into Sprint~1 but was deferred to Sprint~3 where it was retrofitted
across every existing controller. This is discussed further in
Section~\ref{sec:retro}.

\subsection{Scrum Boards (GitHub Projects)}

\subsubsection*{Sprint 1 Boards}

\begin{figure}[H]
  \centering
  \imgOrPlaceholder[width=0.92\textwidth]{assets/github_project_start.png}{Sprint 1 board — start}
  \caption{Sprint~1 GitHub Projects board at sprint start (Backlog populated).}
\end{figure}

\begin{figure}[H]
  \centering
  \imgOrPlaceholder[width=0.92\textwidth]{assets/github_project_mid.png}{Sprint 1 board — mid}
  \caption{Sprint~1 GitHub Projects board mid-sprint (work in progress).}
\end{figure}

\begin{figure}[H]
  \centering
  \imgOrPlaceholder[width=0.92\textwidth]{assets/github_project_end.png}{Sprint 1 board — end}
  \caption{Sprint~1 GitHub Projects board at sprint close (Done column populated).}
\end{figure}

\subsubsection*{Sprint 2 Boards}

\begin{figure}[H]
  \centering
  \imgOrPlaceholder[width=0.92\textwidth]{assets/sprint2_board.png}{Sprint 2 GitHub Projects board}
  \caption{Sprint~2 GitHub Projects board at sprint close --- backlog, ready, in progress, in review, and done columns with all 12 Sprint~2 stories tracked.}
\end{figure}

\subsubsection*{Sprint 3 Boards}

\begin{figure}[H]
  \centering
  \imgOrPlaceholder[width=0.92\textwidth]{../D3/assets/sprintboard1.png}{Sprint 3 board — top}
  \caption{Sprint~3 GitHub Projects board --- top view of the kanban columns.}
\end{figure}

\begin{figure}[H]
  \centering
  \imgOrPlaceholder[width=0.92\textwidth]{../D3/assets/sprintboard2.png}{Sprint 3 board — continued}
  \caption{Sprint~3 GitHub Projects board --- remaining issues and the Done column at sprint close (20~Apr~2026).}
\end{figure}

\subsection{Burn-Down Charts}

\begin{figure}[H]
  \centering
  \imgOrPlaceholder[width=0.78\textwidth]{assets/sprint1_burndown.png}{Sprint 1 Burn-Down Chart}
  \caption{Sprint~1 burn-down chart.}
\end{figure}

\begin{figure}[H]
  \centering
  \imgOrPlaceholder[width=0.78\textwidth]{assets/sprint2_burndown.png}{Sprint 2 Burn-Down Chart}
  \caption{Sprint~2 burn-down chart (planned vs.\ completed story points across 1~--~25~March~2026).}
\end{figure}

\begin{figure}[H]
  \centering
  \imgOrPlaceholder[width=0.78\textwidth]{../D3/assets/burndown-chart.png}{Sprint 3 Burn-Down Chart}
  \caption{Sprint~3 burn-down chart (26~March~--~20~April~2026).}
\end{figure}

% ═══════════════════════════════════════════════════════════════
%  SECTION 13 — WORK DIVISION
% ═══════════════════════════════════════════════════════════════
\section{Work Division (Project-Wide)}

{\sloppy
\subsubsection*{Ibraheem Farooq \normalfont\small\color{muted}(PM / Scrum Master / Developer)}
Sprint planning, sprint tracking, and team coordination across all three
sprints. Ownership of the GitHub Projects board, branch and PR review
discipline, and final packaging of every deliverable. On the implementation
side: end-to-end authentication and RBAC --- tabbed Login/Register UI,
\texttt{Session} singleton, \texttt{PasswordUtil}, \texttt{LoginController},
enrolment-based catalog filtering, and the Access Denied guard on
instructor-only actions; plus announcement file and link attachments
(US\,12--13) including the \texttt{announcement\_attachments} table and chip
rendering in the course stream. Backend support for the course and
announcement modules in Sprint~2.

\subsubsection*{Wajih-Ur-Raza Asif \normalfont\small\color{muted}(Requirements Architect / Tester / Developer)}
Requirements ownership across the project: all 26 user stories, sub-story
breakdowns, and structured specifications, plus consistency verification
between acceptance criteria and actual implemented behaviour. Authored every
deliverable document (D1, D2, D3, and this Final Report). Owned the test
plan, all 25 black-box test cases, white-box branch-coverage analysis on
\texttt{LoginController.handleLogin()} and \texttt{GradingController.handleSave()},
and the bug report. Implemented the Assignments tab and People tab
(US\,25--26): \texttt{Assignment.fetchWithStatusForUser()}, status badge
rendering, the Grade Submissions button, \texttt{Course.fetchEnrolledUsers()},
and roster card rendering. Implemented the full grading pipeline
(US\,19--22): \texttt{grade}/\texttt{feedback} schema additions,
\texttt{Submission} model extensions
(\texttt{fetchByAssignmentId}, \texttt{saveGrade}, \texttt{fetchGradedForStudent},
\texttt{fetchGradebookForCourse}), \texttt{GradingView.fxml} +
\texttt{GradingController}, and \texttt{GradebookView.fxml} +
\texttt{GradebookController} with role-adaptive rendering (student personal
view vs.\ instructor class-wide grid).

\subsubsection*{Ali Aan Khowaja \normalfont\small\color{muted}(Lead Developer / UI Designer)}
Lead implementation of all Sprint~2 features: Announcements, Assignment
Creation, AssignmentDetails with file submission flow and threaded comments,
and the To-Do List with urgency colour-coding. UI design and screen layout
ownership across the entire project --- every FXML view and CSS sheet ---
giving Naawbi a consistent visual language from the login screen through to
the gradebook. Produced implementation screenshots for every deliverable.
\par}

% ═══════════════════════════════════════════════════════════════
%  SECTION 14 — RETROSPECTIVE
% ═══════════════════════════════════════════════════════════════
\section{Retrospective and Lessons Learned}
\label{sec:retro}

\textbf{What went well.} All scoped user stories shipped, including the two
late additions (Assignments tab and People tab) discovered during Sprint~3
planning. The layered MVC architecture established at the start absorbed
every new feature without structural change. The branch-per-feature strategy
kept code review manageable. The idempotent \texttt{seed.sql} with three test
accounts and mixed-state submissions accelerated manual testing across all
three sprints. Single-query gradebook and LEFT-JOIN status derivation
prevented N+1 patterns from creeping in under deadline pressure.

\textbf{What could have been better.} Authentication should have been
Sprint~1 work, not Sprint~3 --- by deferring it, RBAC had to be retrofitted
onto every existing controller late in the project. The grading chain
(database $\rightarrow$ UI $\rightarrow$ gradebook) serialised work that
could have run in parallel if scoped earlier. US\,25--26 (Assignments and
People tabs) should have been identified during Sprint~2 planning when the
assignment infrastructure was fresh, rather than discovered during Sprint~3.

\textbf{What we would do differently.} Ship authentication and RBAC in
Sprint~1 so that every subsequent feature is built session-aware from the
start. Identify course-home navigation tabs (Stream / Assignments / People)
as part of Sprint~2 planning, alongside the assignment creation work. Use
an integration branch for feature chains with linear merge dependencies
instead of waiting on sequential PR reviews.

% ═══════════════════════════════════════════════════════════════
%  SECTION 15 — FINAL SYSTEM SUMMARY
% ═══════════════════════════════════════════════════════════════
\section{Final System Summary}

Naawbi LMS was built across three sprints, each delivering a distinct layer
of the final product.

\begin{center}
\begin{tabular}{@{}p{0.18\linewidth}p{0.13\linewidth}p{0.62\linewidth}@{}}
  \toprule
  \fcol{Sprint} & Deliverable & Core Contribution \\
  \midrule
  \fcol{Sprint 1} & D1 & Course structure --- creation, catalogue, class-code enrolment. \\
  \fcol{Sprint 2} & D2 & Content and submissions --- announcements, assignments, file uploads, status tracking, To-Do list. \\
  \fcol{Sprint 3} & D3 & Identity and evaluation --- authentication, RBAC, announcement attachments, Assignments and People tabs, grading interface, gradebook. \\
  \bottomrule
\end{tabular}
\end{center}

\noindent The final system fulfils its original vision: a course-centric LMS
where instructors manage content and evaluations, and students track progress
and performance, without needing any external tool for the core academic
workflow. \textbf{All 26 planned user stories have been delivered and
verified, all 25 formal black-box test cases pass, and all four bugs found
during testing were fixed before submission.}

\end{document}
